% ===================== 第 4 章（下） 选词填空·阅读·听力 =====================
\section{真题分类汇编：Cloze, Reading and Listening}

\subsection{Banked Cloze (Section A, Reading)}

\begin{engbox}[title=What is tested]
A passage of about 250--300 words with ten blanks. Fifteen words are given; ten must be chosen, and \textbf{each word may be used only once}. The fifteen options always fall into four parts of speech: \textbf{adverbs 2--3, verbs 5--7, adjectives 2--5, nouns 3--5}. Verbs may appear as third-person singular, -ed or -ing forms.
\end{engbox}

\subsubsection*{Real exam items (short excerpts, 2016.12)}

\begin{engbox}
(1) ``Along with this effort, we're working with animal advocates and state and federal lawmakers to \underline{\hspace{2.6cm}} anti-cruelty laws across the country\dots'' \\
--- Answer: \textbf{strengthen} (a bare infinitive after \emph{to}).\\[2pt]
(2) ``The U.S. Sentencing Commission, which \underline{\hspace{2.6cm}} these sentencing guidelines, is revisiting them\dots'' \\
--- Answer: \textbf{creates} (a third-person singular verb in a relative clause).
\end{engbox}

\subsubsection*{Real exam opening (2022.06, Set 1)}

\begin{engbox}
``The city of Bath was founded by the Romans almost two thousand years ago. It has been famous for its \underline{\hspace{2.6cm}} pleasing architecture and healing thermal springs ever since.''\par
Note how the word before \emph{pleasing} must be an \textbf{adverb} modifying the present participle (aesthetically / exceptionally type of word).
\end{engbox}

\subsubsection*{Signals for each part of speech}

\begin{center}
\small
\begin{longtable}{p{3.2cm}p{11.4cm}}
\toprule
\textbf{Gap form} & \textbf{Signal} \\
\midrule
\endfirsthead
\toprule
\textbf{Gap form} & \textbf{Signal} \\
\midrule
\endhead
Noun & after an article / adjective / possessive; between ``and'' or ``or'' and a noun; after a preposition \\
\midrule
Verb & after \emph{to} (bare form), a modal or auxiliary; a missing predicate; a relative clause lacking its verb; parallel with another verb \\
\midrule
Adjective & after \emph{be} / \emph{seem} / \emph{become}; before a noun \\
\midrule
Adverb & between subject and predicate; at the end of a complete sentence; before an adjective or participle \\
\bottomrule
\end{longtable}
\end{center}

\begin{notebox}[title=Five-step method]
1. Read the first sentence (usually gap-free) to catch the topic. \quad 2. Label all fifteen options by part of speech. \quad 3. Predict the part of speech needed for each gap. \quad 4. Compare words of the same part of speech against the context. \quad 5. Read the passage again with your answers in place.
\end{notebox}

\subsection{Long Reading (Section B, Matching)}

\begin{engbox}[title=What is tested]
Ten statements must be matched to the paragraphs of a long article (about 1000 words, 15 paragraphs). The key skill is \textbf{recognizing synonyms and paraphrases}: the statement rarely repeats the wording of the paragraph.
\end{engbox}

\begin{itemize}
  \item Underline the ``anchor words'' in each statement: names, numbers, dates, and unique nouns.
  \item Scan paragraph by paragraph, not statement by statement; mark the paragraph letter beside every promising clue.
  \item Watch for \textbf{synonym swaps}: increase $\to$ rise / grow / soar; important $\to$ significant / crucial; problem $\to$ issue / challenge; reduce $\to$ cut / lower / curb.
  \item Expect one paragraph to be used twice and one or two to be unused.
  \item Do the easy statements first; each confirmed match removes both a statement and a paragraph from the search space.
\end{itemize}

\subsubsection*{Authentic topics from past papers (examples)}

\begin{center}
\small
\begin{longtable}{p{3.6cm}p{11.0cm}}
\toprule
\textbf{Source} & \textbf{Topic reflected in the passage} \\
\midrule
\endfirsthead
\toprule
\textbf{Source} & \textbf{Topic reflected in the passage} \\
\midrule
\endhead
2017.06, Reading & The ``18-hour city'': Houston, Austin, Charlotte and Nashville position themselves as highly competitive in livability, employment and recreational facilities. \\
\midrule
2017.12, Listening B & The great cultural impact of printing was that it facilitated the growth of national languages. \\
\midrule
2019.06, Listening C & A cultural psychologist examines similarities and differences between East Asians and North Americans. \\
\midrule
2017.12, Listening C & Kwanzaa, a cultural festival during which African-Americans celebrate their rich heritage. \\
\midrule
2017.06, Listening C & People can meet their basic needs but lack access to cultural goods, entertainment and recreation. \\
\midrule
2018.12, Listening C & Language diversity must be about history, cultural differences, and mountains or oceans dividing populations. \\
\midrule
2019.12, Listening C & Feeling guilty about the time not spent watching children grow up or travelling. \\
\bottomrule
\end{longtable}
\end{center}

\subsection{Careful Reading (Section C)}

\begin{engbox}[title=What is tested]
Two passages, five multiple-choice questions each, 14.2 points per question --- the highest-value questions in the whole reading section. Question types: main idea, detail, inference, attitude / tone, and vocabulary-in-context.
\end{engbox}

\begin{itemize}
  \item \textbf{Read the questions first}, then the passage, so that details stand out when you meet them.
  \item Locate each answer in the text and underline the sentence; the correct option is a paraphrase, not a copy.
  \item For attitude questions, collect adjectives and verbs that carry judgement.
  \item For inference questions, choose the option that must be true, not the one that merely sounds reasonable.
  \item Past reading passages cluster around: \emph{health, travel, popular science, and social culture}. Prepare background vocabulary in these areas.
\end{itemize}

\subsection{Listening}

\begin{center}
\small
\begin{longtable}{p{3.4cm}p{2.0cm}p{2.4cm}p{6.8cm}}
\toprule
\textbf{Section} & \textbf{Items} & \textbf{Points} & \textbf{Focus} \\
\midrule
\endfirsthead
\toprule
\textbf{Section} & \textbf{Items} & \textbf{Points} & \textbf{Focus} \\
\midrule
\endhead
A. Long conversations & 8 & 7.1 each & Campus life, job interviews, travel plans, shopping and services \\
\midrule
B. Passages & 7 & 7.1 each & Culture, education, science news, social phenomena \\
\midrule
C. Lectures / talks & 10 & 14.2 each & Social science and popular science; the heaviest part \\
\bottomrule
\end{longtable}
\end{center}

\begin{engbox}[title=Exam-room rules that decide your score]
\begin{itemize}
  \item The questions are \textbf{not printed on the paper}; they are read out in the recording. Preview the options during the 25-second intervals.
  \item \textbf{Transfer answers to the answer sheet while listening.} Answer Sheet 1 is collected immediately after the listening section.
  \item If you miss one item, let it go and move to the next question; hesitation costs two questions, not one.
\end{itemize}
\end{engbox}

\begin{engbox}[title=High-frequency listening vocabulary by scenario]
Campus: enrolment, credit, seminar, deadline, scholarship. \quad
Workplace: application, promotion, vacancy, salary, colleague. \quad
Travel: reservation, boarding pass, delay, itinerary, refund. \quad
Science and society: survey, sample, evidence, trend, experiment. \quad
Culture: heritage, festival, ceremony, tradition, ritual.
\end{engbox}
